\chapter{Database}

\section{Database Architecture}

PostgreSQL is the single source of truth for all persistent data. The
application never treats an in-memory or derived value as authoritative
when a real database record could answer the same question --- for
example, worked hours are always recomputed from closed
\code{attendance\_entries} rows, never cached on the \code{Employee}
row itself.

SQLAlchemy 2.0's declarative style (\code{Mapped[...]} type
annotations) defines every table in \file{app/models/*.py}. Each model
file's own docstring documents \emph{why} its constraints exist, not
just what they are --- this documentation quotes many of those
docstrings directly, since they are the authoritative rationale.

\subsection*{Model $\rightarrow$ Migration $\rightarrow$ Table}

Alembic migrations in \file{migrations/versions/} are the actual,
applied schema history; \file{app/models/*.py} is SQLAlchemy's
in-Python \emph{description} of that schema, used for query building
and (via \code{compare\_type=True} in \code{Migrate.init\_app}) for
autogenerating future migrations. The two are cross-checked here rather
than assumed identical, per this project's own documentation
requirement.

Several important constraints in the schema are \textbf{hand-written},
not autogenerated, because SQLAlchemy's Alembic autogeneration cannot
produce PostgreSQL's \code{EXCLUDE ... USING gist} DDL on its own ---
this is called out explicitly in the docstrings of
\longcode{Shift.ex\_shifts\_employee\_no\_overlap},
\longcode{AttendanceEntry.ex\_attendance\_entries\_employee\_no\_overlap},
\longcode{LeaveRequest.ex\_leave\_requests\_employee\_no\_overlap},
\longcode{EmployeePayRate.ex\_employee\_pay\_rates\_employee\_no\_overlap},
and \longcode{OvertimePolicy.ex\_overtime\_policies\_organization\_no\_overlap}.
No discrepancy between the model definitions and the migration scripts
was found during the preparation of this document; both were read
directly for every table below.

\subsection*{Naming Convention}

\file{app/models/base.py} registers a SQLAlchemy naming convention
before any table is defined, so every constraint/index Alembic
autogenerates gets a predictable, greppable name (\code{ix\_},
\code{uq\_}, \code{ck\_}, \code{fk\_}, \code{pk\_} prefixes). This is
why constraint names appear consistently throughout this chapter and in
the service-layer error handling that matches specific constraint names
(e.g.\ \code{uq\_attendance\_entries\_employee\_id\_open}) to translate
a raw database error into a friendly \code{ValidationError}.

\section{Database Table Inventory}

Thirteen tables exist in the schema, verified against
\file{migrations/versions/0001} through \file{0016} and every model
file.

\begin{longtable}{@{}p{2.9cm}p{4.6cm}p{1.6cm}p{4.9cm}@{}}
\toprule
\textbf{Table} & \textbf{Purpose} & \textbf{PK} & \textbf{Important columns} \\
\midrule
\endhead
\code{organizations} & The tenant boundary; every other table is
  scoped to one. & \code{id} & \code{slug} (unique), \code{timezone},
  \code{currency\_code} (CHAR 3), \code{is\_active} \\
\code{departments} & Organizational units employees belong to. & \code{id} &
  \code{organization\_id}, \code{name} (case-insensitive unique per
  org), \code{code}, \code{is\_active} \\
\code{employees} & The HR record: a person who can be scheduled and
  tracked, distinct from their login. & \code{id} & \code{organization\_id},
  \code{department\_id}, \code{employee\_number}, \code{first\_name},
  \code{last\_name}, \code{email} (nullable, unique when present),
  \code{employment\_status}, \code{pay\_type} (\code{'hourly'} only),
  \code{hired\_on}, \code{terminated\_on} \\
\code{users} & A login/authentication account, optionally linked to an
  employee. & \code{id} & \code{organization\_id}, \code{employee\_id}
  (nullable, unique), \code{email} (unique, case-insensitive),
  \code{password\_hash}, \code{role}, \code{is\_active},
  \code{failed\_login\_count}, \code{locked\_until} \\
\code{department\_managers} & Join table: which manager-role users may
  act on which departments. & composite \newline
  (\code{user\_id}, \code{department\_id}) & \code{organization\_id} \\
\code{shifts} & A single planned work period (\emph{scheduling}, not
  attendance). & \code{id} & \code{organization\_id},
  \code{department\_id}, \code{employee\_id} (nullable = unassigned),
  \code{starts\_at}, \code{ends\_at}, \code{business\_date},
  \code{break\_minutes}, \code{status}, \code{published\_at} \\
\code{attendance\_entries} & A single actual work period, recorded via
  clock-in/clock-out. & \code{id} & \code{organization\_id},
  \code{employee\_id}, \code{shift\_id} (nullable),
  \code{started\_at}, \code{ended\_at} (nullable while open),
  \code{business\_date}, \code{break\_minutes}, \code{status},
  \code{source}, \code{edit\_reason} \\
\code{overtime\_policies} & An organization's effective-dated overtime
  configuration. & \code{id} & \code{organization\_id},
  \code{daily\_threshold\_hours}, \code{weekly\_threshold\_hours},
  \code{week\_start\_day}, \code{effective\_from}/\code{\_to} \\
\code{overtime\_tiers} & One multiplier tier belonging to a policy. &
  \code{id} & \code{policy\_id}, \code{scope} (\code{daily}/\code{weekly}),
  \code{tier\_order}, \code{from\_hours}, \code{to\_hours}, \code{multiplier} \\
\code{leave\_types} & A small per-organization catalog of leave
  categories. & \code{id} & \code{organization\_id}, \code{code},
  \code{name}, \code{is\_paid}, \code{requires\_approval},
  \code{blocks\_scheduling}, \code{is\_active} \\
\code{leave\_requests} & An employee's request for time off and its
  approval lifecycle. & \code{id} & \code{organization\_id},
  \code{employee\_id}, \code{leave\_type\_id}, \code{starts\_at}/\code{ends\_at},
  \code{status}, \code{decided\_by\_user\_id}, \code{decided\_at},
  \code{decision\_note} \\
\code{employee\_pay\_rates} & An employee's effective-dated hourly rate
  history. & \code{id} & \code{employee\_id}, \code{organization\_id},
  \code{hourly\_rate} (\code{NUMERIC(10,4)}), \code{effective\_from}/\code{\_to} \\
\code{audit\_logs} & Append-only record of privileged/sensitive actions. &
  \code{id} & \code{organization\_id} (nullable), \code{actor\_user\_id}
  (nullable), \code{action}, \code{entity\_type}, \code{entity\_id},
  \code{changes} (JSONB), \code{ip\_address} (INET), \code{created\_at} \\
\bottomrule
\end{longtable}

\section{Entity Relationship Diagram}

The schema is shown in two parts for readability: the organizational
identity hierarchy, and the workforce/attendance domain tables that
hang off it. Every arrow below corresponds to a real foreign key
(several of them composite, tying a child row's \code{organization\_id}
to its parent's) verified against the model files in \S3.2.

\subsection*{Organizational Identity}

\begin{center}
\resizebox{\textwidth}{!}{%
\begin{tikzpicture}[node distance=1.1cm and 2.2cm]
  \node[flownode, text width=4cm] (org) {\textbf{organizations}\\ \footnotesize id, slug, timezone};
  \node[flownode, text width=4cm, below left=1.4cm and 0.5cm of org] (dept) {\textbf{departments}\\ \footnotesize id, organization\_id, name};
  \node[flownode, text width=4cm, below right=1.4cm and 0.5cm of org] (usr) {\textbf{users}\\ \footnotesize id, organization\_id, employee\_id, role};
  \node[flownode, text width=4cm, below=1.6cm of dept] (emp) {\textbf{employees}\\ \footnotesize id, organization\_id, department\_id};
  \node[flownode, text width=4.4cm, below right=1.6cm and -0.3cm of dept] (dm) {\textbf{department\_managers}\\ \footnotesize user\_id, department\_id, organization\_id};

  \draw[flowarrow] (org) -- node[left,font=\scriptsize]{1:N} (dept);
  \draw[flowarrow] (org) -- node[right,font=\scriptsize]{1:N} (usr);
  \draw[flowarrow] (dept) -- node[left,font=\scriptsize]{1:N} (emp);
  \draw[flowarrow, dashed] (emp) -- node[below,font=\scriptsize]{0..1:1 (nullable, unique)} (usr);
  \draw[flowarrow] (dept) -- node[left,font=\scriptsize]{1:N} (dm);
  \draw[flowarrow] (usr) -- node[right,font=\scriptsize]{1:N} (dm);
\end{tikzpicture}%
}
\end{center}

The dashed \code{employees}$\to$\code{users} edge is deliberately
0..1:1 and \emph{nullable}: an admin or manager account need not have a
linked HR record at all, but a \code{role = 'employee'} account
\emph{must} (enforced by the \code{employee\_role\_requires\_employee\_id}
CHECK constraint, not just application code).

\subsection*{Workforce and Attendance Domain}

\begin{center}
\resizebox{\textwidth}{!}{%
\begin{tikzpicture}[node distance=1.0cm and 1.3cm]
  \node[flownode, text width=3.6cm] (emp) {\textbf{employees}};
  \node[flownode, text width=3.6cm, below left=1.3cm and -0.6cm of emp] (shift) {\textbf{shifts}\\ \footnotesize starts\_at, ends\_at, status};
  \node[flownode, text width=3.6cm, below=1.3cm of emp] (att) {\textbf{attendance\_\allowbreak entries}\\ \footnotesize started\_at, ended\_at, status};
  \node[flownode, text width=3.6cm, below right=1.3cm and -0.6cm of emp] (lr) {\textbf{leave\_requests}\\ \footnotesize starts\_at, ends\_at, status};
  \node[flownode, text width=3.6cm, right=1.4cm of lr] (lt) {\textbf{leave\_types}};
  \node[flownode, text width=3.6cm, left=1.4cm of shift] (pr) {\textbf{employee\_\allowbreak pay\_rates}\\ \footnotesize hourly\_rate, effective\_from/\_to};

  \draw[flowarrow] (emp) -- node[left,font=\scriptsize]{1:N} (shift);
  \draw[flowarrow] (emp) -- node[right,font=\scriptsize]{1:N} (att);
  \draw[flowarrow] (emp) -- node[right,font=\scriptsize]{1:N} (lr);
  \draw[flowarrow] (lt) -- node[above,font=\scriptsize]{1:N} (lr);
  \draw[flowarrow] (emp) -- node[above,font=\scriptsize]{1:N} (pr);
  \draw[flowarrow, dashed] (shift) -- node[below,font=\scriptsize]{0..1:N (nullable)} (att);

  \node[flownode, text width=3.4cm, below=2.6cm of shift.south, anchor=north, xshift=1.6cm] (policy) {\textbf{overtime\_\allowbreak policies}};
  \node[flownode, text width=3.4cm, right=1.4cm of policy] (tier) {\textbf{overtime\_\allowbreak tiers}};
  \node[flownode, text width=3.4cm, right=1.4cm of tier] (audit) {\textbf{audit\_logs}};

  \draw[flowarrow] (policy) -- node[above,font=\scriptsize]{1:N} (tier);
\end{tikzpicture}%
}
\end{center}

\code{overtime\_policies} and \code{audit\_logs} both belong to
\code{organizations} directly (not shown as an edge above to avoid
crossing every other line); \code{audit\_logs.actor\_user\_id} is a
nullable foreign key to \code{users} (nullable because a failed login
against a nonexistent email establishes no user at all). The dashed
\code{shifts}$\to$\code{attendance\_entries} edge is the "matched shift"
link set at clock-in time by \code{attendance.\_match\_shift} --- an
attendance entry can exist with no matched shift (\code{shift\_id IS
NULL}), representing unscheduled work or a clock-in that didn't
unambiguously match exactly one published shift.

\section{Database Relationship Explanation}

\begin{lstlisting}[basicstyle=\ttfamily\small,frame=single]
Organization
  |
  +-- Department
  |      |
  |      +-- Employee
  |             |
  |             +-- User (0..1, the login account)
  |             +-- Shift (planned work)
  |             +-- AttendanceEntry (actual work)
  |             +-- LeaveRequest
  |             +-- EmployeePayRate
  |
  +-- DepartmentManager (User <-> Department, many-to-many)
  +-- OvertimePolicy -- OvertimeTier
  +-- LeaveType
  +-- AuditLog
\end{lstlisting}

\textbf{Organization $\to$ Department $\to$ Employee.} A strict
hierarchy: an employee cannot exist without a department, and a
department cannot exist without an organization. This mirrors the
real-world structure the application models, and every composite
foreign key in the schema (e.g.\ \code{employees.(department\_id,
organization\_id)} $\to$ \code{departments.(id, organization\_id)})
exists specifically to prevent an employee from ever being wired to a
department in a \emph{different} organization than its own.

\textbf{Employee $\leftrightarrow$ User.} Deliberately \emph{not} the
same row. An \code{Employee} is an HR record; a \code{User} is a login
credential. This separation is why an admin or manager can exist with
no employee record at all (they manage the system but are not
necessarily "staff" being scheduled), while every \code{employee}-role
user must have exactly one linked employee record.

\textbf{Department $\leftrightarrow$ User, via DepartmentManager.} A
many-to-many join: one manager can manage several departments, and
(implicitly, though not currently exercised by any route) a department
could have more than one manager. A manager's entire
\code{AccessScope.department\_ids} is sourced from this one table.

\textbf{Employee $\to$ Shift, AttendanceEntry, LeaveRequest,
EmployeePayRate.} Four independent one-to-many relationships, each
representing a different kind of historical record about the same
person. They are kept as separate tables specifically so that, for
example, a pay rate change never has to touch (or invalidate) existing
attendance rows.

\textbf{OvertimePolicy $\to$ OvertimeTier.} The only \code{ON DELETE
CASCADE} relationship in the entire schema (every other foreign key is
\code{RESTRICT}) --- a tier has no independent meaning without its
parent policy, so cascading its deletion is safe; \code{database.md}'s
"avoid accidental cascade deletes" rule is treated as a default to
override only with this kind of genuine ownership relationship, not a
blanket prohibition.

\section{Constraints and Data Integrity}

The schema treats workforce data as financial-grade: money and
attendance figures flow directly into payroll-adjacent calculations, so
integrity is enforced at the database layer wherever practical, not
only in application code.

\subsection*{Overlap Prevention (GiST Exclusion Constraints)}

Five \code{EXCLUDE ... USING gist} constraints are the schema's
signature integrity mechanism, each preventing a specific kind of
double-booking that a simple \code{UNIQUE} constraint cannot express:

\begin{itemize}
  \item \code{ex\_shifts\_employee\_no\_overlap} --- no employee may
        have two overlapping, non-cancelled, assigned shifts.
  \item \code{ex\_attendance\_entries\_employee\_no\_overlap} --- no
        employee may have two overlapping attendance entries (an open
        entry's unbounded range still counts as overlapping anything
        after it starts).
  \item \code{ex\_leave\_requests\_employee\_no\_overlap} --- no
        employee may have two overlapping pending/approved leave
        requests.
  \item \code{ex\_employee\_pay\_rates\_employee\_no\_overlap} --- no
        employee may have two overlapping effective-dated pay rate
        periods.
  \item \code{ex\_overtime\_policies\_organization\_no\_overlap} ---
        an organization may have at most one overtime policy in force
        on any given calendar day.
\end{itemize}

Each of these is paired with a \emph{best-effort} application-level
pre-check (e.g.\ \code{scheduling.\_check\_overlap}) that gives a
friendly error before ever reaching the database --- but the database
constraint is the actual authority under concurrency; every service
function that can trigger one of these catches the resulting
\code{IntegrityError}, inspects \code{error.orig.diag.constraint\_name},
and re-raises a clean \code{ValidationError} instead of letting a raw
500 reach the user.

\subsection*{Partial Unique Indexes}

\code{uq\_attendance\_entries\_employee\_id\_open} (unique on
\code{employee\_id} \code{WHERE ended\_at IS NULL}) is the actual
mechanism preventing a duplicate clock-in: an employee may have at most
one \emph{unresolved} (open or needs-review) attendance entry at a
time, regardless of any application-level check. Similarly,
\code{uq\_employees\_organization\_id\_email} is unique only
\code{WHERE email IS NOT NULL}, so multiple employees with no email on
file never collide with each other.

\subsection*{Key CHECK Constraints}

\begin{itemize}[nosep]
  \item \code{shifts.duration\_max\_24\_hours} /
        \code{attendance\_entries.duration\_max\_24\_hours} --- no
        single shift or attendance entry may span more than 24 hours,
        preventing an arbitrarily long entry from distorting overtime
        or labor cost.
  \item \code{attendance\_entries.status\_open\_matches\_ended\_at\_null}
        --- \code{ended\_at} is \code{NULL} if and only if status is
        \code{open} or \code{needs\_review}; only \code{closed} carries
        an end time.
  \item \code{attendance\_entries.edit\_requires\_edited\_at\_and\_reason}
        --- a correction must always carry who, when, and why together,
        never just one or two of the three.
  \item \code{leave\_requests.decision\_matches\_status} (split into two
        CHECKs) --- carefully designed so that cancelling a
        \emph{previously approved} request preserves who approved it and
        when, distinguishing it from a request cancelled while still
        pending. See the extended rationale in
        \code{app/models/leave\_request.py}.
  \item \code{employees.terminated\_status\_matches\_terminated\_on} ---
        an employee's status is \code{'terminated'} if and only if
        \code{terminated\_on} is set.
  \item \code{users.employee\_role\_requires\_employee\_id} --- a
        \code{role = 'employee'} login must be linked to an employee
        record.
\end{itemize}

\subsection*{Indexes}

Indexes are added only where a real, identified query pattern needs
them (\file{.claude/rules/database.md}: "add indexes only when
justified"), and several were added specifically in response to slow
query shapes discovered during development --- for example,
\longcode{ix\_attendance\_entries\_employee\_id\_business\_date} exists
because \code{working\_hours.worked\_seconds\_for\_day}\allowbreak/\code{\_week}
are called in tight per-day loops by cost and report calculations, and
neither pre-existing index matched that predicate shape.

\subsection*{Cascade Behavior}

Every foreign key in the schema is \code{ON DELETE RESTRICT}
\textbf{except} \code{overtime\_tiers.policy\_id}, which is
\code{ON DELETE CASCADE} (a tier cannot outlive its policy) and
\code{department\_managers}'s foreign key to \code{users}, which is
also \code{ON DELETE CASCADE} (a manager assignment has no meaning once
the user account itself is gone). \code{RESTRICT} everywhere else means
the schema will refuse to delete, for example, a department that still
has employees --- deactivation (\code{is\_active = false}), not
deletion, is how the application retires a department or an employee
record, preserving historical schedules/attendance/cost data.
